%=============================================================================
\section{Simulation Workflow}
%=============================================================================

You have a discretized controller and a mathematical model of your plant.
Before burning the code into flash and risking a \$200 motor, why not test it
on your laptop first? This chapter shows you how to build simulations that
catch 90\% of bugs before they become hardware damage.

We will start with the case for simulation, then walk through three levels of
testing fidelity---Model-in-the-Loop (MIL), Software-in-the-Loop (SIL), and
Hardware-in-the-Loop (HIL)---before showing how to inject realism into your
models. A copy-paste Python template and a balance-robot case study wrap things
up.

%-----------------------------------------------------------------------------
\subsection{Why Simulate Before Deploying}
%-----------------------------------------------------------------------------

Hardware debugging is expensive in every dimension that matters during
competition season. Broken gears cost money. Burned motor windings cost
money \textit{and} a two-week shipping delay. Even when nothing breaks, every
hour spent staring at an oscilloscope is an hour not spent tuning your
auto-aim or refining your chassis path planner.

Simulation cannot replace hardware testing, but it dramatically reduces the
number of hardware iterations you need:

\textbf{What simulation CAN tell you:}
\begin{itemize}
    \item \textbf{Stability.} If the closed-loop poles are in the right half of
          the $z$-plane (or equivalently outside the unit circle), your gains
          are wrong. Simulation catches this instantly; hardware catches it by
          destroying a gearbox.
    \item \textbf{Approximate performance.} Rise time, overshoot, settling time
          will be within 20--40\% of reality for a decent model.
    \item \textbf{Gain sensitivity.} Sweep a parameter $\pm 50\%$ and see if
          the controller still works. If simulation shows marginal stability
          at $+30\%$ inertia, your real robot will oscillate the first time
          someone bolts on a heavier armor plate.
\end{itemize}

\textbf{What simulation CANNOT tell you:}
\begin{itemize}
    \item Real friction (stiction, Coulomb, temperature-dependent).
    \item Backlash in geartrains---unless you explicitly model it.
    \item Exact sensor noise characteristics (quantization, drift, outliers).
    \item True communication latency (CAN bus arbitration, DMA jitter).
\end{itemize}

The simulation-reality gap is unavoidable. Accept it. The goal is not a
perfect simulation; the goal is a simulation good enough that your first
hardware test produces ``it basically works, just needs gain tweaking'' instead
of ``the motor spun backwards and ripped the wiring harness off.''

%-----------------------------------------------------------------------------
\subsection{Model-in-the-Loop (MIL)}
%-----------------------------------------------------------------------------

MIL is the simplest and most common simulation approach: you write both the
plant model and the controller in the same language (Python, MATLAB, or even
C), run them in a loop, and look at the results.

\subsubsection{Building a DC Motor + PID Simulation from Scratch}

We will use the simplified DC motor model from Chapter~3. Neglecting
inductance (valid when the electrical time constant $\tau_e = L/R$ is much
smaller than the control period $T_s$), the mechanical dynamics are:
\begin{equation}
J \frac{d\omega}{dt} = \frac{K_t}{R}\bigl(V - K_e \omega\bigr) - b\,\omega
\end{equation}
where $J$ is rotor inertia, $K_t$ the torque constant, $K_e$ the back-EMF
constant, $R$ the winding resistance, $b$ the viscous friction coefficient,
$V$ the applied voltage, and $\omega$ the angular velocity.

Discretize with forward Euler ($T_s$ = sampling period):
\begin{equation}
\omega[n+1] = \omega[n] + T_s \cdot
    \frac{1}{J}\left[\frac{K_t}{R}\bigl(V[n] - K_e\,\omega[n]\bigr)
    - b\,\omega[n]\right]
\end{equation}

Pair this with the discrete PID from Chapter~4:
\begin{equation}
u[n] = K_p\,e[n] + K_i\,T_s \sum_{k=0}^{n} e[k]
       + \frac{K_d}{T_s}\bigl(e[n] - e[n-1]\bigr)
\end{equation}

Here is a complete Python simulation:

\begin{verbatim}
import numpy as np
import matplotlib.pyplot as plt

# Plant parameters (small brushed DC motor)
J, b, Kt, Ke, R = 0.01, 0.1, 0.01, 0.01, 1.0
Ts, N = 0.001, 2000          # 1 kHz, 2 seconds
omega = np.zeros(N)
# PID gains
Kp, Ki, Kd = 5.0, 10.0, 0.05
e_int, e_prev = 0.0, 0.0
omega_ref = 100.0             # rad/s step reference

for n in range(N - 1):
    e = omega_ref - omega[n]
    e_int += e * Ts
    e_dot = (e - e_prev) / Ts
    V = Kp * e + Ki * e_int + Kd * e_dot
    V = np.clip(V, -24.0, 24.0)          # actuator saturation
    domega = ((Kt / R) * (V - Ke * omega[n]) - b * omega[n]) / J
    omega[n + 1] = omega[n] + Ts * domega
    e_prev = e

t = np.arange(N) * Ts
plt.plot(t, omega); plt.axhline(omega_ref, ls='--', c='r')
plt.xlabel('Time (s)'); plt.ylabel('omega (rad/s)')
plt.title('DC Motor PID Step Response'); plt.grid(); plt.show()
\end{verbatim}

Twenty-five lines and you have a working simulation. Run it, and you will see
a step response that overshoots slightly and settles in about 0.5~s. If it
diverges, your gains are wrong---better to learn that on a laptop than on a
\$200 motor.

\subsubsection{Comparing Simulation with Theory}

After collecting the step response, measure the rise time $t_r$, overshoot
$M_p$, and settling time $t_s$ directly from the simulation data. Then compare
with the second-order approximation formulas from Chapter~3:
\begin{align}
M_p &= e^{-\pi\zeta / \sqrt{1 - \zeta^2}} \\
t_s &\approx \frac{4}{\zeta \omega_n}
\end{align}

If the numbers match within 10--15\%, your simulation is consistent with the
linearized model. If they do not match:
\begin{enumerate}
    \item Your simulation might have a bug (check units, check signs).
    \item The closed-loop system is not well-approximated by second order
          (dominant poles are not a clean conjugate pair).
    \item Actuator saturation is clipping the response and invalidating the
          linear analysis.
\end{enumerate}

The second-order formulas are approximations, not ground truth. When in doubt,
trust the simulation---it includes nonlinearities (like saturation) that the
formulas ignore.

%-----------------------------------------------------------------------------
\subsection{Software-in-the-Loop (SIL)}
%-----------------------------------------------------------------------------

MIL is great for design exploration, but it has a blind spot: your competition
robot does not run Python. It runs C on an STM32, with fixed-point arithmetic
(or at least 32-bit float), specific integer types, and ISR-based timing.

The idea behind SIL is simple: \textbf{compile your actual embedded C
controller code and run it against the Python plant model.} This catches an
entire class of bugs that MIL misses:
\begin{itemize}
    \item Fixed-point overflow (your integral term wraps around at 65535).
    \item Integer division truncation (\texttt{int16\_t} gains lose precision).
    \item Variable type mismatches (\texttt{uint8\_t} holding a value that
          should be \texttt{int32\_t}).
    \item ISR timing assumptions that do not hold.
\end{itemize}

The implementation is straightforward. Compile your C controller into a shared
library and call it from Python using \texttt{ctypes}:

\begin{verbatim}
import ctypes

# Compile: gcc -shared -o pid.so -fPIC pid.c
lib = ctypes.CDLL('./pid.so')
lib.pid_update.argtypes = [ctypes.c_float, ctypes.c_float]
lib.pid_update.restype  = ctypes.c_float
lib.pid_init()

for n in range(N - 1):
    V = lib.pid_update(ctypes.c_float(omega_ref),
                       ctypes.c_float(omega[n]))
    # ... plant update same as before ...
\end{verbatim}

Now the PID running in your simulation is byte-for-byte identical to the one
that will run on your STM32. If the integral term overflows in C, it overflows
in your simulation too---and you catch it before competition day.

\textbf{When to skip SIL:} If your controller is pure Python or MATLAB (e.g.,
you are doing rapid prototyping and have not written the embedded code yet),
MIL is sufficient. SIL only adds value when you have actual C/C++ code to
test.

%-----------------------------------------------------------------------------
\subsection{Hardware-in-the-Loop (HIL)}
%-----------------------------------------------------------------------------

HIL takes things one step further: your \textbf{real MCU} runs your
\textbf{real code}, but the plant is still simulated on a PC. The PC and MCU
communicate via UART, SPI, or CAN---the PC receives the control output from
the MCU, simulates the plant dynamics for one time step, and sends back
simulated sensor readings.

\textbf{The setup:}
\begin{enumerate}
    \item MCU boots and starts its control loop normally.
    \item Instead of reading a real encoder, it reads a value sent by the PC
          over UART.
    \item It computes the PID output and sends the PWM duty cycle (or voltage
          command) back to the PC.
    \item The PC updates the plant model and sends the new ``sensor'' value.
    \item Repeat at the control rate.
\end{enumerate}

\textbf{When HIL is worth the effort:}
\begin{itemize}
    \item Safety-critical systems where a controller bug could injure someone.
    \item Very expensive actuators (industrial servos, hydraulic valves).
    \item Systems where the real plant is not available yet (waiting for
          custom-machined parts).
\end{itemize}

\textbf{When HIL is overkill:} Most competition scenarios. If your motor costs
\$30 and your gearbox is not going to explode from a gain that is 2x too high,
go straight from MIL to real hardware. The time you spend setting up HIL
communication protocols would be better spent tuning on the real robot.

\textbf{Latency warning:} HIL at 1~kHz means you have $\leq 1$~ms round-trip
for UART communication + plant simulation + sensor encoding. At 115200~baud,
sending 8 bytes takes about 0.7~ms. Use higher baud rates (921600 or USB FS)
or reduce the control rate for HIL testing.

%-----------------------------------------------------------------------------
\subsection{Making Simulations Realistic}
%-----------------------------------------------------------------------------

A perfect simulation is useless for controller validation. If your PID only
works when the plant model is exact and the sensor is noiseless, it will fail
on real hardware. You need to inject the same kinds of imperfections that
reality will throw at you.

\subsubsection{Sensor Noise}

Add white Gaussian noise with variance matched to your sensor's datasheet (or
measured Allan variance). For a typical MEMS gyroscope with noise density
$0.01~\text{deg/s}/\sqrt{\text{Hz}}$ sampled at 1~kHz, the discrete noise
standard deviation is $\sigma = 0.01 \times \sqrt{1000} \approx
0.316$~deg/s.

\subsubsection{Disturbances}

\begin{itemize}
    \item \textbf{Step disturbance:} apply a sudden load torque at a random
          time to test disturbance rejection.
    \item \textbf{Sinusoidal disturbance:} inject $d(t) = A\sin(2\pi f t)$ at
          a known frequency to test frequency-domain rejection.
\end{itemize}

\subsubsection{Friction}

The simplest useful friction model combines viscous and Coulomb friction:
\begin{equation}
\tau_{\text{friction}} = b\,\omega + \tau_c \cdot \operatorname{sign}(\omega)
\end{equation}
where $b$ is the viscous coefficient (already in our motor model) and
$\tau_c$ is the Coulomb friction torque. This model captures the deadband
around zero velocity that makes low-speed tracking difficult.

\subsubsection{Communication Delay}

Delay the control signal by $N_d$ samples to model CAN bus latency, sensor
processing time, or computation delay:
$V_{\text{applied}}[n] = V_{\text{computed}}[n - N_d]$.

\subsubsection{Actuator Saturation}

Clamp the output to $\pm V_{\max}$ (or $\pm \text{PWM}_{\max}$). This is
critical---linear analysis assumes unbounded control effort, but your battery
is 24~V and your ESC clips at 100\% duty.

Here is how to add noise, friction, and saturation to the simulation loop:

\begin{verbatim}
tau_c = 0.005       # Coulomb friction torque (N*m)
sigma = 0.3         # sensor noise std dev (rad/s)
N_delay = 3         # 3-sample control delay
V_buf = [0.0] * (N_delay + 1)

for n in range(N - 1):
    omega_meas = omega[n] + np.random.randn() * sigma    # noisy sensor
    e = omega_ref - omega_meas
    # ... PID computation ...
    V_buf.append(np.clip(V_cmd, -24, 24))                # saturation
    V = V_buf[-(N_delay + 1)]                             # delayed output
    friction = b * omega[n] + tau_c * np.sign(omega[n])   # friction
    domega = ((Kt / R) * (V - Ke * omega[n]) - friction) / J
    omega[n + 1] = omega[n] + Ts * domega
\end{verbatim}

\textbf{Key insight:} A controller that works in a perfect simulation but
fails when you add noise and friction is \textit{not} a good controller. If
your gains only work in paradise, they will not survive competition day.

%-----------------------------------------------------------------------------
\subsection{Quick Python Simulation Template}
%-----------------------------------------------------------------------------

Here is a complete, copy-paste-ready template. Change the plant parameters to
match your robot and you have a working simulation in two minutes.

\begin{verbatim}
import numpy as np, matplotlib.pyplot as plt

def simulate(plant_fn, Kp, Ki, Kd, ref, Ts=0.001, dur=2.0,
             noise_std=0.0, V_max=24.0):
    N = int(dur / Ts)
    y, u = np.zeros(N), np.zeros(N)
    e_int, e_prev, state = 0.0, 0.0, 0.0
    for n in range(N - 1):
        y_meas = y[n] + np.random.randn() * noise_std
        e = ref - y_meas
        e_int += e * Ts
        u_cmd = Kp * e + Ki * e_int + Kd * (e - e_prev) / Ts
        u[n] = np.clip(u_cmd, -V_max, V_max)
        y[n+1], state = plant_fn(y[n], u[n], state, Ts)
        e_prev = e
    t = np.arange(N) * Ts
    fig, (a1, a2) = plt.subplots(2, 1, sharex=True)
    a1.plot(t, y); a1.axhline(ref, ls='--', c='r'); a1.set_ylabel('Output')
    a2.plot(t, u); a2.set_ylabel('Control'); a2.set_xlabel('Time (s)')
    plt.tight_layout(); plt.show()
    return t, y, u
\end{verbatim}

Define a plant function for your specific system. For the DC motor:

\begin{verbatim}
def dc_motor(omega, V, _, Ts,
             J=0.01, b=0.1, Kt=0.01, Ke=0.01, R=1.0):
    domega = ((Kt/R)*(V - Ke*omega) - b*omega) / J
    return omega + Ts * domega, None

t, y, u = simulate(dc_motor, Kp=5, Ki=10, Kd=0.05,
                    ref=100, noise_std=0.3)
\end{verbatim}

Copy this, change the plant function parameters to match your robot, and you
are running.

%-----------------------------------------------------------------------------
\vspace{1em}
\begin{mdframed}[linewidth=1pt, innertopmargin=10pt, innerbottommargin=10pt]
\textbf{Case Study: Tuning LQR in Simulation Before Deploying to a Balance
Robot}

\vspace{0.3em}
A two-wheeled balance robot is a classic inverted pendulum on wheels. The state
vector is $\mathbf{x} = [\theta,\;\dot\theta,\;x,\;\dot x]^T$---tilt angle
and angular rate of the body, plus position and velocity of the wheels. The
linearized continuous-time dynamics around the upright equilibrium are:
\begin{equation}
\dot{\mathbf{x}} = A\mathbf{x} + Bu, \qquad
A = \begin{bmatrix}
0 & 1 & 0 & 0 \\
\frac{g}{l} & -\frac{b_\theta}{J_\theta} & 0 & 0 \\
0 & 0 & 0 & 1 \\
-\frac{g m}{M} & 0 & 0 & -\frac{b_x}{M}
\end{bmatrix}, \quad
B = \begin{bmatrix} 0 \\ -\frac{1}{J_\theta l} \\ 0 \\ \frac{1}{M} \end{bmatrix}
\end{equation}
where $g$ is gravity, $l$ the distance from wheel axle to center of mass, $m$
the pendulum mass, $M$ the total mass, and $b_\theta$, $b_x$ are friction
coefficients.

\textbf{Step 1: Discretize.} Using ZOH at $T_s = 5$~ms: $A_d = e^{A T_s}$,
$B_d = A^{-1}(A_d - I)B$. In Python, use
\texttt{scipy.signal.cont2discrete((A, B, C, D), Ts, method='zoh')}.

\textbf{Step 2: Design LQR.} Choose cost matrices $Q$ and $R$ from Chapter~6.
Start conservative:
\[
Q = \operatorname{diag}(100,\; 1,\; 10,\; 1), \qquad R = 1
\]
This penalizes tilt angle 100x more than angular rate, and position 10x more
than velocity. Solve the discrete algebraic Riccati equation to get
$K = [K_\theta,\; K_{\dot\theta},\; K_x,\; K_{\dot x}]$.

\textbf{Step 3: Simulate.} Run the closed-loop system:
$\mathbf{x}[n+1] = A_d \mathbf{x}[n] + B_d(-K\mathbf{x}[n])$ with initial
condition $\theta_0 = 5^\circ$.

Now the fun part: \textbf{sweep $Q$ and $R$} to explore the design space.
\begin{itemize}
    \item Increase $Q_{11}$ (tilt penalty): faster tilt recovery, more
          aggressive control, larger wheel accelerations.
    \item Increase $Q_{33}$ (position penalty): robot drifts less but may
          sacrifice tilt performance.
    \item Increase $R$: gentler control, slower response, lower current draw.
\end{itemize}

After 20 minutes of sweeping in simulation, the team settled on
$Q = \operatorname{diag}(500,\; 5,\; 20,\; 1)$, $R = 0.5$, which gave a
predicted settling time of 0.8~s with less than $2^\circ$ overshoot.

\textbf{Step 4: Deploy.} Flash the gains into the STM32, power on the robot,
give it a push. Hardware result: settling time 1.1~s, overshoot about
$2.5^\circ$. The 30\% gap is mostly unmodeled friction in the gearbox and
slightly higher center of mass than assumed. But---and this is the important
part---\textbf{the gains worked on the first try}. No motor burned, no gears
stripped, no wasted afternoon. One round of minor gain tweaking on hardware
brought the settling time down to 0.9~s.

\textbf{Lesson:} Simulation will never be exact. But it gets you into the
right neighborhood, and that is worth more than any amount of trial-and-error
on hardware.
\end{mdframed}
